\section{Discussion and conclusion}
\label{sec:discussion}

Executing the returned program changes what a mesh-to-CAD search appears to achieve. Stronger search raises reported mesh IoU on every step we measured, but it also produces programs that do not build, and for the rule package and the full ladder the executed-program gain is measurably smaller than the mesh gain. Whole-part section proposals are the extension whose gain holds up best under execution. The same operations, offered as tools, raise every one of six open-weight models well above what it achieves by writing CadQuery alone; the best agent's mean differs from the pipeline's by \mGemmaVsPipeDiff{} on the same parts. Whether the LLM adds anything over a script that uses the same tools is the question our data leave open. Answering it is the natural next experiment, because CAD-agent comparisons that omit such a policy cannot separate the tools from the model.

\paragraph{Limitations.}
(i) Each search arm and each agent condition ran once per part, except three Gemma runs; our intervals capture part sampling, not run-to-run variation. (ii) The agent study uses \nAgentParts{} parts, which leaves intervals wide; the policy comparison in particular is underpowered. (iii) The cohort is hard-tier benchA/benchE, dominated by axis-aligned 2.5D parts: agents rarely called \texttt{fit\_revolve}, so revolve-heavy parts are untested. (iv) We audit one optimizer family. The mesh--program gap we report is a property of these searches, and other reconstructors need their own replay audit. (v) The agent's revolve tool includes a profile-split fix that the frozen pipeline arms lack, and the section-proposal recipe originated from a control designed after seeing agent results on test parts, although it was frozen before the fresh cohort was drawn. (vi) All conditions select answers with IoU against the target mesh; this is the task's input, but results do not transfer to settings without that feedback.
